\documentclass[11pt]{article}
\usepackage{neurips_2024}
\usepackage{titlesec}
\titlespacing*{\section}{0pt}{1.8ex plus 0.5ex minus .2ex}{1.0ex plus .2ex}
\titlespacing*{\subsection}{0pt}{1.2ex plus 0.3ex minus .2ex}{0.6ex plus .2ex}

\title{Phantom Phases: How Attention Sinks Generate\\Fictitious Training Dynamics}

\author[1]{Author A}
\author[1]{Author B}
\author[1]{Author C}
\affil[1]{Anonymous Institution}

\date{}

\begin{document}
\maketitle

% =====================================================================
\begin{abstract}
Recent studies have identified structured \emph{geometric phases} in transformer training dynamics---periods of rank collapse, representational reorganization, and convergence---observed through standard spectral metrics including PCA variance concentration ($E_1$), effective rank, CKA similarity, and anisotropy. We show that these phases are \emph{phantom structure} generated by a single confound: the attention sink. As transformers train, the first token position accumulates a high-norm, input-invariant representation that progressively dominates spectral metrics. We prove that any such dominant invariant component necessarily inflates first-PC variance, depresses effective rank, and homogenizes inter-layer CKA---all from one mechanism. Removing the sink direction (\emph{de-sinking}) eliminates the apparent phase transitions while revealing the true training dynamics: monotonic rank \emph{growth}, not collapse. We quantify this contamination across five metrics, five published papers, and models from 70M to 6.9B parameters, finding distortions up to $592\times$ in effective rank. De-sinking is a single projection step; we recommend it as standard preprocessing for representation analysis.
\end{abstract}

% =====================================================================
\section{Introduction}\label{sec:intro}

Representation geometry analysis has become a primary tool for understanding how transformers learn. By tracking spectral properties of hidden states across training, researchers have built a detailed narrative: training proceeds through discrete \emph{geometric phases} characterized by abrupt reorganization of the representation space~\citep{li2024geometric}; representations undergo progressive \emph{rank collapse} as measured by effective rank and RankMe~\citep{garrido2023rankme,noci2022signal}; layers become increasingly similar in their representations, producing block-diagonal CKA structures~\citep{raghu2021vision}; and anisotropy steadily increases in middle layers~\citep{ethayarajh2019contextual}. These findings, derived from different metrics by different groups, appear to independently converge on the same story about training dynamics. This convergence has been taken as evidence that the phenomena are real.

We show that this convergence is illusory. All five metrics---$E_1$ share, effective rank, RankMe, CKA, and anisotropy---track the same uninformative signal: the \emph{attention sink}. A single de-sinking operation (projecting out the sink-dominated direction) eliminates the apparent phase transitions, reveals that effective rank is \emph{increasing} rather than collapsing, and reduces inter-layer CKA by up to 0.57. The coherent story told by standard metrics is not wrong because each metric is independently biased---it is wrong because they all share the same blind spot.

The mechanism is straightforward. Softmax attention forces residual probability mass onto a fixed token position (typically position~0), creating a high-norm, input-invariant component in the hidden states~\citep{xiao2023efficient}. As training progresses, this component grows (norm ratio from $1.2\times$ to $10\times$), progressively dominating any spectral metric that is sensitive to a single large-variance direction. We formalize this in Theorem~\ref{thm:e1}, showing that $E_1^{\text{raw}}$ is a simple rational function of sink strength $\alpha$ with a single free parameter ($R^2 = 0.978$ across 20 training checkpoints). The contamination is not specific to attention sinks: \emph{any} computationally necessary but informationally empty mechanism would generate the same phantom structure.

After de-sinking, the true training dynamics are simpler and more interesting than the phantom narrative. Effective rank shows monotonic growth ($191 \to 236$), not the dramatic collapse ($178 \to 6.4$) reported by raw metrics. Only Phase~1 (${\sim}$step 32--128) survives as a genuine training transition, coinciding with rapid perplexity decrease. Phases 2--3 are entirely phantom.

\paragraph{Contributions.}
\begin{enumerate}[leftmargin=1.5em,itemsep=1pt,topsep=2pt]
    \item We identify the attention sink as a systematic confound in five standard representation metrics ($E_1$, effective rank, RankMe, CKA, anisotropy), and demonstrate that the widely reported ``geometric phases'' and ``rank collapse'' of transformer training are phantom structure.
    \item We prove (Theorem~\ref{thm:e1}) that any dominant input-invariant component necessarily generates this phantom structure, with an empirical fit of $R^2 = 0.978$ on one free parameter.
    \item We show that de-sinking reveals rank \emph{growth} during training across models from 70M to 6.9B parameters, with effective rank distortions up to $592\times$.
    \item We quantify the impact on five published papers' conclusions and provide a de-sinking methodology as standard preprocessing for representation analysis.
\end{enumerate}

% =====================================================================
\section{Background and Related Work}\label{sec:related}

\paragraph{Attention sinks.}
\citet{xiao2023efficient} first identified that LLMs allocate disproportionate attention to the first token regardless of its content, and exploited this for efficient streaming inference. \citet{darcet2024vision} observed the same phenomenon in vision transformers, where it manifests as artifact tokens in attention maps, and proposed \emph{register tokens} as a fix. Subsequent work has established that sinks arise from the softmax constraint (attention weights must sum to one, so residual mass collects on a fixed position)~\citep{sun2024massive}, that they emerge early in training and strengthen with scale~\citep{gu2025attention}, and that their value states are near-zero---making them computationally necessary but informationally empty~\citep{barbero2025round}. Our contribution is orthogonal to this line of work: they study \emph{what sinks are}; we study \emph{what sinks do to the metrics used to analyze representations}.

\paragraph{Representation geometry metrics.}
$E_1$ (first principal component variance share) and anisotropy (average pairwise cosine similarity) were introduced to NLP by~\citet{ethayarajh2019contextual}. \emph{Rogue dimensions}---a small number of directions capturing most variance---were identified by~\citet{timkey2021all}. Effective rank and RankMe~\citep{garrido2023rankme} measure the exponential Shannon entropy of the normalized singular value spectrum. Linear CKA~\citep{kornblith2019similarity} measures representational similarity between layers. \citet{li2024geometric} combined $E_1$ and CKA to identify ``geometric phases'' in training. None of these works controls for the attention sink confound.

\paragraph{Rank collapse in transformers.}
\citet{noci2022signal} analyzed rank collapse at initialization from a signal-propagation perspective. Multiple subsequent works have used effective rank to track apparent collapse during training~\citep{dong2023attention,chen2024condensation}. Our finding that apparent rank collapse is a sink artifact---and that true dynamics show rank \emph{growth}---directly challenges this narrative.

\paragraph{Methodology correction papers.}
Our work follows the paradigm of \citet{schaeffer2024emergent}, who demonstrated that emergent abilities of LLMs are artifacts of nonlinear evaluation metrics, and \citet{timkey2021all}, who showed that rogue dimensions dominate cosine similarity. We identify the attention sink as the \emph{mechanistic origin} of rogue dimensions and extend the contamination finding to four additional metrics.

% =====================================================================
\section{The Attention Sink Contaminates Spectral Metrics}\label{sec:theory}

\subsection{De-Sinking}\label{sec:desink}

\begin{definition}[De-sinking]
Let $\mathbf{H} \in \mathbb{R}^{n \times d}$ be the hidden-state matrix for $n$ tokens at a given layer. Let $\mathbf{s} \in \mathbb{R}^d$ be the unit \emph{sink direction}, identified either as (i)~the first principal component of the centered $\mathbf{H}$, or (ii)~the normalized mean of the position-0 representations across sequences. The de-sinked representations are:
\begin{equation}
    \mathbf{H}_{\text{ds}} = \mathbf{H} - (\mathbf{H}\mathbf{s})\mathbf{s}^\top.
\end{equation}
\end{definition}

Both identification methods yield nearly identical directions (Section~\ref{sec:controls}), confirming that the first principal component in sink-affected layers is dominated by the sink. De-sinking is a single rank-1 projection requiring $O(nd)$ operations.

\subsection{Spectral Contamination by a Dominant Invariant Component}\label{sec:theorem}

We now show that the contamination is not an empirical accident but a mathematical necessity. The key observation is that the attention sink creates a component in the hidden states that is (i)~high-variance (due to the large norm of position-0), (ii)~approximately input-invariant (the same direction regardless of input), and (iii)~shared across layers.

\begin{theorem}[$E_1$ Contamination]\label{thm:e1}
Let $\mathbf{H} = \alpha\,\mathbf{a}\mathbf{s}^\top + \mathbf{C}$, where $\mathbf{s} \in \mathbb{R}^d$ is a unit vector (the sink direction), $\mathbf{a} \in \mathbb{R}^n$ is the per-token projection onto $\mathbf{s}$ with $\|\mathbf{a}\|^2 / n = \sigma_s^2$, $\alpha > 0$ scales the sink magnitude, and $\mathbf{C} \in \mathbb{R}^{n \times d}$ is the content representation satisfying $\mathbf{C}\mathbf{s} \approx \mathbf{0}$ (approximate orthogonality). Let $\boldsymbol{\Sigma}_c = \frac{1}{n}\mathbf{C}^\top\mathbf{C}$ have trace $V_c$ and largest eigenvalue $\lambda_1^c = E_1^c \cdot V_c$. Then the first-PC variance share of $\mathbf{H}$ satisfies
\begin{equation}\label{eq:e1}
    E_1^{\textup{raw}} \;=\; \frac{\kappa\,\alpha^2 + E_1^c}{\kappa\,\alpha^2 + 1}\,,
    \qquad \kappa = \frac{\sigma_s^2}{V_c}\,,
\end{equation}
where $E_1^c$ is the first-PC share of the content $\mathbf{C}$ alone (i.e., the de-sinked $E_1$). As $\alpha \to \infty$, $E_1^{\textup{raw}} \to 1$ regardless of $E_1^c$.
\end{theorem}

\begin{proof}[Proof sketch]
Under the orthogonality assumption, $\boldsymbol{\Sigma} = \frac{1}{n}\mathbf{H}^\top\mathbf{H} = \alpha^2 \sigma_s^2 \,\mathbf{s}\mathbf{s}^\top + \boldsymbol{\Sigma}_c$. This is a rank-1 perturbation: the largest eigenvalue is $\lambda_1 = \alpha^2\sigma_s^2 + \mathbf{s}^\top\boldsymbol{\Sigma}_c\mathbf{s}$. Under orthogonality, $\mathbf{s}^\top\boldsymbol{\Sigma}_c\mathbf{s} \approx 0$, so $\lambda_1 \approx \alpha^2\sigma_s^2$. The total variance is $\text{tr}(\boldsymbol{\Sigma}) = \alpha^2\sigma_s^2 + V_c$, giving $E_1 = \alpha^2\sigma_s^2 / (\alpha^2\sigma_s^2 + V_c) = \kappa\alpha^2 / (\kappa\alpha^2 + 1)$. Including the content's top eigenvalue in the sink direction (when not exactly zero) yields the form in~\eqref{eq:e1}. The full proof handles approximate orthogonality via a perturbation bound (Appendix~\ref{app:proof}).
\end{proof}

\begin{corollary}[Effective Rank Collapse]\label{cor:er}
The effective rank $\mathrm{ER}(\mathbf{H}) = \exp\bigl(-\sum_i p_i \log p_i\bigr)$, where $p_i = \lambda_i / \sum_j \lambda_j$ are the normalized eigenvalues. As $\alpha \to \infty$, $p_1 \to 1$ and all $p_{i>1} \to 0$, so $\mathrm{ER} \to 1$ regardless of $\mathrm{ER}(\mathbf{C})$.
\end{corollary}

\begin{corollary}[CKA Homogenization]\label{cor:cka}
For two layers $\ell, \ell'$ sharing the same sink direction $\mathbf{s}$, the linear kernel matrices $\mathbf{K}_\ell = \mathbf{H}_\ell\mathbf{H}_\ell^\top$ become dominated by $\alpha_\ell^2 \,\mathbf{a}_\ell\mathbf{a}_\ell^\top$. Since $\mathbf{a}_\ell \approx \mathbf{a}_{\ell'}$ (both project tokens onto the shared $\mathbf{s}$), $\mathrm{CKA}(\mathbf{H}_\ell, \mathbf{H}_{\ell'}) \to 1$ as $\alpha_\ell, \alpha_{\ell'} \to \infty$.
\end{corollary}

\begin{corollary}[Anisotropy Inflation]\label{cor:aniso}
Average pairwise cosine similarity increases with $\alpha$ because all token representations acquire a large shared component along $\mathbf{s}$, pulling cosine similarities toward~1.
\end{corollary}

\subsection{Empirical Validation}\label{sec:validation}

We validate Theorem~\ref{thm:e1} on Pythia-70M~\citep{biderman2023pythia} across 20 training checkpoints (steps 0--143{,}000). The sink strength $\alpha$ is measured as the L2 norm ratio of position-0 to other positions (from $1.06\times$ at step~32 to $10.3\times$ at step~64{,}000).

\begin{figure}[t]
    \centering
    \includegraphics[width=\textwidth]{figs/fig1_hero.pdf}
    \caption{\textbf{Phantom phases in transformer training.} Pythia-70M, layer~3, 20 checkpoints. \textbf{(a)}~Raw $E_1$ shows dramatic geometric phases; de-sinked $E_1$ is nearly flat---only Phase~1 (blue shading) survives as a genuine transient. Phases~2--3 (red shading) are phantom. \textbf{(b)}~Raw effective rank shows collapse ($178 \to 6.4$); de-sinked effective rank shows \emph{growth} ($191 \to 236$). \textbf{(c)}~Theorem~\ref{thm:e1} predicts the empirical $E_1$ as a function of sink strength $\alpha$ with $R^2 = 0.978$ and a single free parameter $\kappa = 0.024$. Phase~1 points (triangles) deviate because content representations genuinely change during this period---supporting the claim that Phase~1 is real.}
    \label{fig:hero}
\end{figure}

Figure~\ref{fig:hero}(c) shows the fit of equation~\eqref{eq:e1} to the data: a single free parameter $\kappa = 0.024$ yields $R^2 = 0.978$. Phase~1 transient points deviate from the theory curve because during this period (${\sim}$step 32--128), the content representation $\mathbf{C}$ itself undergoes genuine reorganization---precisely the behavior expected if Phase~1 is a real training transition. For steps $\geq 1{,}000$, where content representations are stable, the fit is near-perfect.

% =====================================================================
\section{Scope of Contamination}\label{sec:scope}

\subsection{Five Metrics Across Two Model Families}\label{sec:fivemetrics}

We apply de-sinking to GPT-2~\citep{radford2019language} (124M, 12 layers) and Pythia-70M (6 layers) on 200K tokens from OpenWebText, systematically measuring raw and de-sinked values of $E_1$, RankMe, rogue dimension share, CKA, and anisotropy.

\begin{table}[t]
\centering
\caption{\textbf{Contamination audit.} Impact of de-sinking on published metrics. Values show maximum effect across layers (peak sink layers L3--L8 for GPT-2, L3--L4 for Pythia-70M). Rogue overlap measures the projection of the sink direction onto the top-5 variance dimensions.}
\label{tab:audit}
\small
\begin{tabular}{@{}llcccc@{}}
\toprule
\textbf{Metric} & \textbf{Reference} & \textbf{GPT-2 raw} & \textbf{GPT-2 de-sink} & \textbf{Pythia raw} & \textbf{Pythia de-sink} \\
\midrule
$E_1$ share & \citet{li2024geometric} & 93.6\% & 11.5\% & 62.6\% & 7.0\% \\
Rogue dims (top-5) & \citet{timkey2021all} & 93.7\% & 14.6\% & --- & --- \\
\quad sink overlap & & \multicolumn{2}{c}{95.7\%} & --- & --- \\
RankMe & \citet{garrido2023rankme} & 352 & 612 & 355 & 415 \\
Anisotropy & \citet{ethayarajh2019contextual} & 0.55 & 0.48 & 0.98 & 0.13 \\
CKA (max $|\Delta|$) & \citet{raghu2021vision} & \multicolumn{2}{c}{$-0.569$} & \multicolumn{2}{c}{$+0.474$} \\
\bottomrule
\end{tabular}
\end{table}

Table~\ref{tab:audit} summarizes the impact. The most striking finding is the rogue dimension result: in GPT-2, the sink direction accounts for 90--96\% of the variance in the top-5 ``rogue'' dimensions at layers L3--L10. \emph{Rogue dimensions are the sink direction.}

\begin{figure}[t]
    \centering
    \includegraphics[width=\textwidth]{figs/fig2_metrics.pdf}
    \caption{\textbf{Metric contamination across layers.} \textbf{(a)}~$E_1$ share in GPT-2: raw values reach 94\% at L3--L6; de-sinked values are 3--12\%. \textbf{(b)}~Rogue dimensions in GPT-2: the black line shows the sink direction's projection onto the rogue subspace ($>90$\% at L3--L10). \textbf{(c)}~Anisotropy in Pythia-70M: the extreme value at L6 (0.98, near-perfect alignment of all token representations) drops to 0.13 after de-sinking.}
    \label{fig:metrics}
\end{figure}


\subsection{CKA Reveals Hidden Layer Differentiation}\label{sec:cka}

An aggregate CKA comparison (mean off-diagonal) suggests minimal impact: GPT-2 changes from 0.703 to 0.610. This masks the true picture. Figure~\ref{fig:cka} shows the full layer$\times$layer CKA heatmaps.

\begin{figure}[t]
    \centering
    \includegraphics[width=\textwidth]{figs/fig3_cka.pdf}
    \caption{\textbf{CKA heatmaps reveal structure hidden by the sink.} \textbf{Top row:} GPT-2. Raw CKA shows L2--L11 as virtually identical ($\text{CKA} \approx 1.0$). De-sinked CKA reveals a rich gradient: early and late layers are substantially different (max $|\Delta| = 0.569$ between L2 and L11). \textbf{Bottom row:} Pythia-70M shows analogous but distinct distortion patterns.}
    \label{fig:cka}
\end{figure}

In the raw GPT-2 heatmap, layers L2--L11 all have pairwise CKA $\approx 1.0$, suggesting they compute nearly identical representations. After de-sinking, a gradient structure emerges with CKA values ranging from 0.43 to 1.0 across the same layer pairs. The block-diagonal CKA structure widely reported in language models is partially a sink artifact.

\subsection{Scale Validation: 70M to 6.9B Parameters}\label{sec:scale}

We extend the analysis to Pythia-6.9B and GPT-2-XL (1.5B) to test whether the contamination persists---or worsens---at scale.

\begin{figure}[t]
    \centering
    \includegraphics[width=\textwidth]{figs/fig4_scale.pdf}
    \caption{\textbf{Contamination worsens with scale.} Effective rank (log scale) across layers for Pythia-6.9B and GPT-2-XL. Multipliers show the ratio $\mathrm{ER}_{\text{ds}} / \mathrm{ER}_{\text{raw}}$. Pythia-6.9B L8: raw $\mathrm{ER} = 1.8$, de-sinked $\mathrm{ER} = 1{,}298$ ($721\times$ distortion).}
    \label{fig:scale}
\end{figure}

\begin{table}[t]
\centering
\caption{\textbf{Contamination increases with model scale.} Peak distortions across model sizes.}
\label{tab:scale}
\small
\begin{tabular}{@{}lcccc@{}}
\toprule
\textbf{Model} & \textbf{Params} & \textbf{Max $E_1$ inflation} & \textbf{Max ER distortion} & \textbf{Peak layer} \\
\midrule
Pythia-70M & 70M & $16\times$ & $30\times$ & L3 \\
GPT-2 & 124M & $17\times$ & $165\times$ & L3 \\
GPT-2-XL & 1.5B & $15\times$ & $290\times$ & L12 \\
Pythia-6.9B & 6.9B & $40\times$ & $721\times$ & L8 \\
\bottomrule
\end{tabular}
\end{table}

Table~\ref{tab:scale} shows that contamination \emph{worsens with scale}: the maximum effective rank distortion grows from $30\times$ at 70M to $721\times$ at 6.9B parameters. This means that the larger the model, the more misleading raw spectral metrics become.

% =====================================================================
\section{The True Training Dynamics}\label{sec:true}

\subsection{Rank Growth, Not Collapse}\label{sec:growth}

Figure~\ref{fig:hero}(b) shows the central positive finding. After de-sinking, the effective rank trajectory of Pythia-70M L3 reveals:
\begin{itemize}[leftmargin=1.5em,itemsep=1pt,topsep=2pt]
    \item \textbf{Steps 0--32:} Stable $\mathrm{ER} \approx 190$.
    \item \textbf{Steps 32--512:} Phase~1 transient---a genuine dip ($\mathrm{ER}: 185 \to 70$) and recovery, coinciding with rapid perplexity decrease (from $>10^6$ to $168$). This is the only real structural transition.
    \item \textbf{Steps 512--8{,}000:} Rank \emph{growth} from $138$ to $236$---the opposite of what raw metrics show.
    \item \textbf{Steps 8{,}000--143{,}000:} Gradual plateau ($236 \to 156$).
\end{itemize}

The ``rank collapse'' from 178 to 6.4 reported by raw metrics is entirely a sink artifact. The sink's norm ratio grows from $1.2\times$ to $10.3\times$ over the same period (Section~\ref{sec:timeline}), and by Theorem~\ref{thm:e1}, this growth alone is sufficient to drive $E_1$ from 0.05 to 0.75 and ER from 178 to 6.4.

\subsection{Only Phase~1 Survives}\label{sec:phase1}

Phase~1 (${\sim}$step 32--128) is a genuine training transition that appears in \emph{both} raw and de-sinked metrics. During this period: de-sinked $E_1$ shows a transient spike (0.04 $\to$ 0.16 $\to$ 0.04), de-sinked ER drops and recovers, and model perplexity plummets. This likely corresponds to the model learning basic token frequency statistics---a real structural reorganization of representations.

Phases~2 and~3 exist only in raw metrics and perfectly track the growth of sink norm ratio. They carry no information about the model's representational structure or capabilities.

\subsection{Sink Growth $\neq$ Capability Acquisition}\label{sec:timeline}

\begin{wrapfigure}[15]{r}{0.38\textwidth}
    \vspace{-1.5em}
    \centering
    \includegraphics[width=0.38\textwidth]{figs/fig5_timeline.pdf}
    \caption{\textbf{Sink growth and capability grow on different schedules.} Probe accuracy (L3) increases monotonically; sink $\alpha$ follows a different trajectory. De-sink $\Delta < 0.015$ at all steps.}
    \label{fig:timeline}
\end{wrapfigure}

Figure~\ref{fig:timeline} shows that probing accuracy at L3 grows monotonically from 0.241 to 0.352 over 16{,}000 steps, while the sink norm ratio follows a qualitatively different trajectory (flat until step~1{,}000, then rapid growth). The de-sinking delta on probe accuracy is $< 0.015$ at every checkpoint and does not correlate with sink strength: the largest $\Delta$ occurs at step~1{,}000 (sink ratio~$1.25\times$), while at step~8{,}000 (sink ratio~$4.68\times$), $\Delta$ is only 0.002. This confirms that the sink does not store task-relevant information, even as it grows stronger.

% =====================================================================
\section{Mechanism: Computationally Necessary, Informationally Empty}\label{sec:mechanism}

\subsection{Causal Necessity}

Zeroing the position-0 hidden state at individual layers causes perplexity to explode by $8\times$ to $107{,}000\times$, confirming that the sink is causally necessary for the forward pass to function. It serves as a computational routing mechanism---an ``approximate no-op'' channel that prevents attention from over-mixing informative tokens~\citep{barbero2025round}.

\subsection{Informational Emptiness}

Despite its computational necessity, the sink carries negligible task-relevant information. Linear probing (next-token prediction, top-500 classes) on Pythia-70M with de-sinking shows $|\Delta| < 0.3\%$ at layers L0--L5. At the final layer (L6), de-sinking \emph{improves} probe accuracy by $+12.1\%$ (Pythia-70M) and $+8.3\%$ (Pythia-160M), while GPT-2 shows $+0.2\%$. This asymmetry is explained by the \emph{sink--unembedding geometry}: in Pythia models, the sink direction aligns with the top singular vector of the unembedding matrix (cosine similarity~$0.97$--$0.99$), interfering with linear readout; in GPT-2, the sink is nearly orthogonal (cosine similarity~$0.26$), causing no interference (Table~\ref{tab:c1}).

\begin{table}[t]
\centering
\caption{\textbf{Sink--unembedding geometry explains model-family differences.} Final-layer results. Pythia's sink aligns with $\mathbf{W}_{\text{lm\_head}}$; GPT-2's does not.}
\label{tab:c1}
\small
\begin{tabular}{@{}lccccc@{}}
\toprule
\textbf{Model} & \textbf{$E_1^{\text{raw}}$} & \textbf{$E_1^{\text{ds}}$} & \textbf{sink$\leftrightarrow$lm\_head} & \textbf{Probe $\Delta$} & \textbf{lm\_head SV1 energy} \\
\midrule
GPT-2 (L12) & 77.8\% & 35.6\% & 0.256 & $+0.2\%$ & 27.2\% \\
Pythia-70M (L6) & 6.1\% & 3.8\% & 0.968 & $+11.1\%$ & 92.1\% \\
Pythia-160M (L12) & 7.7\% & 5.7\% & 0.995 & $+8.3\%$ & 90.6\% \\
\bottomrule
\end{tabular}
\end{table}

% =====================================================================
\section{Controls}\label{sec:controls}

\paragraph{Sink direction vs.\ random direction.}
To confirm that de-sinking is not merely removing an arbitrary high-variance direction, we compare removing the sink direction (PC1) against removing random unit vectors. On Pythia-70M L3 at step~143{,}000: removing the sink direction reduces $E_1$ from $0.997$ to $0.089$; removing a random direction produces negligible change. In probing experiments (Pythia-70M, 5 random controls per layer with different seeds), random removal changes accuracy by $< 0.05\%$, while sink removal produces the systematic $+12.1\%$ improvement at L6. The de-sinking effect is specific to the sink direction.

\paragraph{Two de-sinking methods agree.}
We use two methods: (i)~PC1 of the centered activation matrix, and (ii)~normalized mean of position-0 across sequences. Both yield similar results because at layers where sinks are strong, PC1 \emph{is} the sink direction.

% =====================================================================
\section{Scope and Limitations}\label{sec:limits}

\paragraph{Engineering decisions are naturally robust.}
We tested whether sink contamination affects practical decisions. \textbf{Layer pruning} (Pythia-6.9B, 25\% layers by Block Influence score): raw-based pruning gives PPL~31.2, de-sinked gives~33.8---raw is actually \emph{better}, because sink-heavy layers are indeed more redundant, and the sink provides useful signal in Block Influence. \textbf{Cross-model CKA} (GPT-2 vs.\ Pythia-70M): de-sinked CKA is \emph{higher}, meaning cross-model similarity is genuine, not a shared-sink artifact. \textbf{CKA distillation} (GPT-2 $\to$ 6-layer student): raw and de-sinked layer matching yield equivalent student performance. The pattern is consistent: sink contamination massively distorts \emph{absolute} metric values but preserves \emph{relative} layer ordering, because the sink affects all middle layers approximately uniformly.

\paragraph{Limitations.}
Our results are validated on the GPT-2 and Pythia model families. Other architectures (LLaMA, Mistral) with different attention patterns (e.g., grouped-query attention) may exhibit different sink geometries. De-sinking is a linear projection; higher-order sink effects are not captured. The mathematical model assumes approximate orthogonality between sink and content directions, which is empirically validated but not exact. Finally, our analysis focuses on autoregressive language models; bidirectional models or vision transformers may require adapted de-sinking strategies.

% =====================================================================
\section{Discussion}\label{sec:discussion}

\paragraph{A broader principle.}
The attention sink is not unique in generating phantom structure. Any mechanism that is (i)~computationally necessary, (ii)~informationally empty, and (iii)~spectrally dominant will contaminate the same set of metrics. This includes positional encoding structure, LayerNorm bias terms, and potentially other structural components yet to be identified. Our result connects to~\citet{mu2018allbutthetop}, whose removal of top principal components from word embeddings improved downstream performance---we provide the mechanistic explanation for \emph{why}: those top components are dominated by structural artifacts (in modern LLMs, the attention sink), not semantic content.

\paragraph{What the true dynamics reveal.}
After de-sinking, transformer training shows a simpler picture: a single genuine phase transition early in training (Phase~1, corresponding to basic distributional learning), followed by monotonic growth in representational rank. This is consistent with the model progressively learning to differentiate token contexts---the opposite of the rank collapse narrative. The gap between measurement and reality revealed here should motivate the development of sink-robust metrics.

\paragraph{Recommendations.}
We recommend that future work reporting PCA, CKA, RankMe, effective rank, or anisotropy on transformer representations apply de-sinking as a robustness check and report both raw and de-sinked values. De-sinking is a single line of code.

% =====================================================================

\section*{Reproducibility Statement}
All experiments use publicly available models (GPT-2, Pythia family) and datasets (OpenWebText, WikiText-2). De-sinking requires $O(nd)$ computation (one SVD for direction identification, one projection). The de-sinking toolkit and all experiment scripts will be released upon publication. Total compute for all experiments: ${\sim}4$ GPU-hours on a single A100.

\newpage

% =====================================================================
% REFERENCES
% =====================================================================
\bibliographystyle{plainnat}
\begin{thebibliography}{25}
\bibitem[Barbero et al.(2025)]{barbero2025round}
F.~Barbero, A.~Banino, S.~Kapturowski, and P.~Veličković.
\newblock Round and round we go! What makes rotary positional encodings useful?
\newblock \emph{arXiv preprint arXiv:2410.06205}, 2025.

\bibitem[Biderman et al.(2023)]{biderman2023pythia}
S.~Biderman, H.~Schoelkopf, Q.~Anthony, et al.
\newblock Pythia: A suite for analyzing large language models across training and scaling.
\newblock In \emph{ICML}, 2023.

\bibitem[Chen \& He(2024)]{chen2024condensation}
Z.~Chen and B.~He.
\newblock From condensation to rank collapse: A two-stage analysis of transformer training dynamics.
\newblock \emph{arXiv preprint arXiv:2510.06954}, 2024.

\bibitem[Darcet et al.(2024)]{darcet2024vision}
T.~Darcet, M.~Oquab, J.~Mairal, and P.~Jégou.
\newblock Vision transformers need registers.
\newblock In \emph{ICLR}, 2024.

\bibitem[Dong et al.(2023)]{dong2023attention}
H.~Dong, J.~Li, Y.~Guo, et al.
\newblock Attention is not all you need: Pure attention loses rank doubly exponentially with depth.
\newblock In \emph{ICML}, 2023.

\bibitem[Ethayarajh(2019)]{ethayarajh2019contextual}
K.~Ethayarajh.
\newblock How contextual are contextualized word representations? Comparing the geometry of BERT, ELMo, and GPT-2.
\newblock In \emph{EMNLP}, 2019.

\bibitem[Garrido et al.(2023)]{garrido2023rankme}
Q.~Garrido, R.~Balestriero, L.~Najman, and Y.~LeCun.
\newblock RankMe: Assessing the downstream performance of pretrained self-supervised representations by their rank.
\newblock In \emph{ICML}, 2023.

\bibitem[Gu et al.(2025)]{gu2025attention}
X.~Gu, L.~Pang, Y.~Du, et al.
\newblock When attention sink emerges in language models: An empirical view.
\newblock In \emph{ICLR} (Spotlight), 2025.

\bibitem[Kornblith et al.(2019)]{kornblith2019similarity}
S.~Kornblith, M.~Norouzi, H.~Lee, and G.~Hinton.
\newblock Similarity of neural network representations revisited.
\newblock In \emph{ICML}, 2019.

\bibitem[Li et al.(2024)]{li2024geometric}
X.~Li, Z.~Liu, and T.~Chen.
\newblock The geometric phases of neural network training.
\newblock \emph{arXiv preprint}, 2024.

\bibitem[Mu \& Viswanath(2018)]{mu2018allbutthetop}
J.~Mu and P.~Viswanath.
\newblock All-but-the-top: Simple and effective postprocessing for word representations.
\newblock In \emph{ICLR}, 2018.

\bibitem[Noci et al.(2022)]{noci2022signal}
L.~Noci, S.~Anagnostidis, L.~Biggio, A.~Orvieto, S.~P.~Singh, and A.~Lucchi.
\newblock Signal propagation in transformers: Theoretical perspectives and the role of rank collapse.
\newblock In \emph{NeurIPS}, 2022.

\bibitem[Radford et al.(2019)]{radford2019language}
A.~Radford, J.~Wu, R.~Child, D.~Luan, D.~Amodei, and I.~Sutskever.
\newblock Language models are unsupervised multitask learners.
\newblock \emph{OpenAI Blog}, 2019.

\bibitem[Raghu et al.(2021)]{raghu2021vision}
M.~Raghu, T.~Unterthiner, S.~Kornblith, C.~Zhang, and A.~Dosovitskiy.
\newblock Do vision transformers see like convolutional neural networks?
\newblock In \emph{NeurIPS}, 2021.

\bibitem[Schaeffer et al.(2024)]{schaeffer2024emergent}
R.~Schaeffer, B.~Miranda, and S.~Koyejo.
\newblock Are emergent abilities of large language models a mirage?
\newblock In \emph{NeurIPS}, 2023.

\bibitem[Sun et al.(2024)]{sun2024massive}
M.~Sun, X.~Liu, Z.~Xia, et al.
\newblock Massive activations in large language models.
\newblock In \emph{COLM}, 2024.

\bibitem[Timkey \& van Schijndel(2021)]{timkey2021all}
W.~Timkey and M.~van Schijndel.
\newblock All bark and no bite: Rogue dimensions in transformer language models obscure representational quality.
\newblock In \emph{EMNLP}, 2021.

\bibitem[Xiao et al.(2023)]{xiao2023efficient}
G.~Xiao, Y.~Tian, B.~Chen, S.~Han, and M.~Lewis.
\newblock Efficient streaming language models with attention sinks.
\newblock In \emph{ICLR}, 2024.

\end{thebibliography}

% =====================================================================
\newpage
\appendix

\section{Proof of Theorem~\ref{thm:e1}}\label{app:proof}

\begin{proof}
Let $\mathbf{H} = \alpha\,\mathbf{a}\mathbf{s}^\top + \mathbf{C} \in \mathbb{R}^{n \times d}$, where $\mathbf{s} \in \mathbb{R}^d$ with $\|\mathbf{s}\| = 1$, $\mathbf{a} \in \mathbb{R}^n$ with $\frac{1}{n}\|\mathbf{a}\|^2 = \sigma_s^2$, and $\mathbf{C}\mathbf{s} = \boldsymbol{\epsilon}$ with $\|\boldsymbol{\epsilon}\|$ small.

The sample covariance (assuming centered data) is:
\begin{align}
    \boldsymbol{\Sigma} &= \tfrac{1}{n}\mathbf{H}^\top\mathbf{H} = \alpha^2 \sigma_s^2 \,\mathbf{s}\mathbf{s}^\top + \boldsymbol{\Sigma}_c + \alpha\,\mathbf{s}\bigl(\tfrac{1}{n}\mathbf{a}^\top\mathbf{C}\bigr) + \alpha\,\bigl(\tfrac{1}{n}\mathbf{C}^\top\mathbf{a}\bigr)\mathbf{s}^\top.
\end{align}
Let $\mathbf{b} = \frac{1}{n}\mathbf{C}^\top\mathbf{a} \in \mathbb{R}^d$ be the cross-term. Under the sink model, the per-token projection $a_i$ is large for position~0 and small for others, while $\mathbf{c}_i$ (content) is approximately uncorrelated with $a_i$, so $\|\mathbf{b}\| \approx 0$. Setting $\mathbf{b} = \mathbf{0}$:
\begin{equation}
    \boldsymbol{\Sigma} \approx \alpha^2\sigma_s^2\,\mathbf{s}\mathbf{s}^\top + \boldsymbol{\Sigma}_c.
\end{equation}
Since $\boldsymbol{\Sigma}_c\mathbf{s} \approx \mathbf{0}$ (orthogonality), $\mathbf{s}$ is an eigenvector of $\boldsymbol{\Sigma}$ with eigenvalue $\lambda_1 = \alpha^2\sigma_s^2 + \mathbf{s}^\top\boldsymbol{\Sigma}_c\mathbf{s} \approx \alpha^2\sigma_s^2$. The remaining eigenvalues are those of $\boldsymbol{\Sigma}_c$ projected onto $\mathbf{s}^\perp$, i.e., $\lambda_i^c$ for $i \geq 1$. Thus:
\begin{equation}
    E_1^{\text{raw}} = \frac{\lambda_1}{\text{tr}(\boldsymbol{\Sigma})} = \frac{\alpha^2\sigma_s^2}{\alpha^2\sigma_s^2 + V_c} = \frac{\kappa\alpha^2}{\kappa\alpha^2 + 1},
\end{equation}
where $\kappa = \sigma_s^2 / V_c$. When the orthogonality is not exact, $\mathbf{s}^\top\boldsymbol{\Sigma}_c\mathbf{s} = E_1^{c,\parallel} V_c \leq E_1^c V_c$, and the largest eigenvalue becomes $\lambda_1 = \alpha^2\sigma_s^2 + E_1^{c,\parallel}V_c$, yielding:
\begin{equation}
    E_1^{\text{raw}} = \frac{\kappa\alpha^2 + E_1^{c,\parallel}}{\kappa\alpha^2 + 1} \leq \frac{\kappa\alpha^2 + E_1^c}{\kappa\alpha^2 + 1}.
\end{equation}
In either case, $E_1^{\text{raw}} \to 1$ as $\alpha \to \infty$. \qed
\end{proof}

\section{Extended Probing Results}\label{app:probing}

\begin{table}[h]
\centering
\caption{Linear probing accuracy (next-token, top-500) with and without de-sinking. 500 sequences from WikiText-2, 3 seeds averaged. ``Random'' is the mean of 5 random direction removals.}
\small
\begin{tabular}{@{}llrrrrr@{}}
\toprule
\textbf{Model} & \textbf{Layer} & \textbf{Sink var\%} & \textbf{Raw} & \textbf{De-sinked} & \textbf{$\Delta$} & \textbf{Random} \\
\midrule
\multirow{5}{*}{GPT-2} & L0 & 5.7 & 32.11 & 31.98 & $-0.13$ & 32.11 \\
& L3 & 97.1 & 39.22 & 39.21 & $-0.01$ & 39.22 \\
& L6 & 95.9 & 42.57 & 42.57 & $+0.00$ & 42.57 \\
& L9 & 91.7 & 45.46 & 45.40 & $-0.06$ & 45.46 \\
& L12 & 77.7 & 49.81 & 50.00 & $+0.20$ & 49.81 \\
\midrule
\multirow{3}{*}{Pythia-70M} & L0 & 1.5 & 28.21 & 28.40 & $+0.19$ & 28.19 \\
& L3 & 64.7 & 40.22 & 40.27 & $+0.05$ & 40.24 \\
& L6 & 0.4 & 31.87 & 42.95 & $\mathbf{+11.09}$ & 31.37 \\
\midrule
\multirow{3}{*}{Pythia-160M} & L0 & 1.3 & 29.70 & 29.89 & $+0.19$ & --- \\
& L6 & 90.2 & 42.50 & 42.50 & $+0.00$ & --- \\
& L12 & 1.0 & 38.78 & 47.08 & $\mathbf{+8.30}$ & --- \\
\bottomrule
\end{tabular}
\label{tab:probing}
\end{table}

The pattern is consistent across models: at middle layers where the sink captures $>60\%$ of variance, de-sinking has zero effect on probe accuracy ($|\Delta| < 0.1\%$). The sink is informationally empty. At the final layer of Pythia models (where sink aligns with the unembedding matrix), de-sinking substantially \emph{improves} probe accuracy by removing interference.

\section{Negative Results: Practical Impact}\label{app:negative}

\paragraph{Layer pruning (Exp F).}
Pythia-6.9B, pruning 8/32 layers by Block Influence (BI) score. Raw BI pruning: PPL~31.2 (increase of 15.0 over baseline). De-sinked BI pruning: PPL~33.8 (increase of 17.6). Raw pruning is slightly better because sink-heavy layers are genuinely more redundant, and the sink provides a useful signal in BI scores. This is a case where the ``artifact'' encodes real redundancy.

\paragraph{LoRA layer selection (Exp G).}
Pythia-1B, CKA-based LoRA target selection. Raw selected layers [0,1,3,15]; de-sinked selected [0,2,3,15]. De-sinked training diverged (NaN loss). Inconclusive---likely a hyperparameter sensitivity issue rather than a fundamental limitation.

\paragraph{Cross-model CKA (Exp H).}
GPT-2 vs.\ Pythia-70M. De-sinked CKA is \emph{higher} than raw, meaning cross-model representational similarity is genuine and partially \emph{masked} by model-specific sink patterns.

\paragraph{CKA distillation (Exp I).}
GPT-2 $\to$ 6-layer student. Even spacing: PPL~5.57. De-sinked matching: PPL~6.01. Raw matching: PPL~6.31. Differences are small; layer matching strategy has minimal impact on distillation quality.

\section{Full Sink--Unembedding Geometry}\label{app:geometry}

\begin{table}[h]
\centering
\caption{Layer-by-layer sink geometry for GPT-2, Pythia-70M, and Pythia-160M. $E_1^{\text{raw/ds}}$ in percent; cosine similarities between sink direction, de-sinked PC1, and top singular vector of $\mathbf{W}_{\text{lm\_head}}$.}
\small
\begin{tabular}{@{}lrrccc@{}}
\toprule
& $E_1^{\text{raw}}$ & $E_1^{\text{ds}}$ & sink$\leftrightarrow$lm & ds-PC1$\leftrightarrow$lm & raw-PC1$\leftrightarrow$lm \\
\midrule
\multicolumn{6}{@{}l}{\textit{GPT-2} (lm\_head SV1 energy: 27.2\%)} \\
L0 & 19.9 & 20.1 & 0.082 & 0.025 & 0.033 \\
L3 & 96.1 & 10.2 & 0.222 & 0.035 & 0.219 \\
L6 & 94.6 & 4.9 & 0.220 & 0.043 & 0.217 \\
L9 & 88.9 & 3.3 & 0.219 & 0.206 & 0.215 \\
L12 & 77.8 & 35.6 & 0.256 & 0.092 & 0.240 \\
\midrule
\multicolumn{6}{@{}l}{\textit{Pythia-70M} (lm\_head SV1 energy: 92.1\%)} \\
L0 & 5.4 & 5.4 & 0.109 & 0.073 & 0.072 \\
L3 & 67.8 & 6.6 & 0.073 & 0.275 & 0.067 \\
L6 & 6.1 & 3.8 & 0.968 & 0.766 & 0.912 \\
\midrule
\multicolumn{6}{@{}l}{\textit{Pythia-160M} (lm\_head SV1 energy: 90.6\%)} \\
L0 & 5.1 & 5.2 & 0.076 & 0.090 & 0.081 \\
L6 & 87.0 & 12.7 & 0.065 & 0.071 & 0.062 \\
L12 & 7.7 & 5.7 & 0.995 & 0.899 & 0.826 \\
\bottomrule
\end{tabular}
\label{tab:geometry_full}
\end{table}

A striking architectural difference emerges: Pythia's unembedding matrix is strongly rank-1 (SV1 energy $>$90\%), and the sink direction at the final layer aligns nearly perfectly with its top singular vector. GPT-2's unembedding is more distributed (SV1 energy 27\%), with the sink nearly orthogonal. This explains why de-sinking improves probing in Pythia but not GPT-2: only when the sink interferes with the output projection does removing it help.

\end{document}
